\chapter{Loss of Female Libido}
\index{Loss of Female Libido}

\section*{Definition and Overview}
Loss of libido, or low sexual desire, is a decrease in the interest in, or desire for, sexual activity. It is the most common sexual problem reported by women, and it is important to recognise that sexual desire varies naturally between women and across a woman's life. Low desire becomes a concern when it causes personal distress or relationship problems. The medical term hypoactive sexual desire disorder is used when low desire is persistent, causes distress, and is not better explained by another condition. The causes are usually a mixture of physical, psychological, and relationship factors, and treatment addresses the underlying drivers rather than offering a simple fix. This chapter explains why desire changes and what can help.

\section*{Epidemiology}
Low sexual desire affects a large proportion of women at some point. Studies suggest that roughly one in three women experiences a loss of interest in sex at some time, and that around 1 in 10 has persistent low desire that causes distress. Desire typically fluctuates with life stage, being affected by childbearing, menopause, stress, and relationship duration. Women report low desire more often than men, but this reflects both biological and social factors. Despite its frequency, it is greatly under-treated, because women often do not raise it and doctors often do not ask.

\section*{Aetiology and Pathophysiology}
\begin{itemize}
    \item \textbf{Physical factors:} hormonal changes around menstruation, pregnancy, breastfeeding, and menopause; thyroid disorders; chronic illness; pain; fatigue; and the side effects of medicines including antidepressants, blood pressure drugs, and some contraceptives
    \item \textbf{Psychological factors:} stress, anxiety, depression, low self-esteem, body image concerns, and a history of trauma or abuse
    \item \textbf{Relationship factors:} poor communication, unresolved conflict, loss of intimacy, and differing expectations between partners
    \item \textbf{Sexual problems:} pain during sex (dyspareunia), vaginal dryness, or difficulty reaching orgasm can reduce desire because sex is anticipated as unpleasant
    \item \textbf{Lifestyle:} alcohol, drug use, poor sleep, and a demanding work or caring load
    \item \textbf{Hormonal mechanisms:} oestrogen and testosterone influence desire, and low levels, especially after menopause or removal of the ovaries, can reduce libido
    \item \textbf{Responsive desire:} many women experience desire in response to intimacy rather than spontaneously, so context and relationship quality matter greatly
\end{itemize}

\section*{Clinical Features}
\begin{itemize}
    \item Reduced or absent interest in sexual activity, fantasies, or initiation of sex
    \item Fewer sexual thoughts and less response to physical or emotional cues
    \item Avoidance of sex or participating only to please a partner
    \item Distress, frustration, guilt, or anxiety about the change
    \item Effects on the relationship: distance, conflict, or reduced intimacy
    \item Associated symptoms may point to causes: pain with sex, dryness, low mood, fatigue, or menopausal symptoms
    \item The pattern matters: lifelong versus recent onset, general versus situation-specific, and whether desire returns in new relationships
\end{itemize}

\section*{Differential Diagnosis}
A range of conditions can cause or contribute to low desire.
\begin{itemize}
    \item Depression and anxiety disorders
    \item Hypoactive sexual desire disorder, the specific diagnosis when desire is low and distressing
    \item Sexual pain disorders, including dyspareunia and vaginismus
    \item Menopause and genitourinary syndrome of menopause
    \item Thyroid disease, diabetes, and other medical conditions
    \item Medication effects, particularly antidepressants and hormonal contraceptives
    \item Substance use, including alcohol
    \item Relationship distress and sexual trauma
\end{itemize}

\section*{When to Seek Medical Care}
Women should feel able to raise low desire with a doctor or a sexual health or women's health service. Medical review is important to exclude treatable physical causes and to identify contributing medicines.

\textbf{Contact your local emergency services for:}
\begin{itemize}
    \item Sudden loss of desire with neurological symptoms such as weakness or headache, which is rare but requires urgent assessment
    \item Any symptoms suggesting stroke, heart attack, or another emergency
    \item Suicidal thoughts or feelings, which can accompany depression underlying low desire
\end{itemize}

\section*{Diagnosis and Investigations}
\begin{itemize}
    \item \textbf{Detailed history:} onset, duration, context, relationship history, medicines, alcohol, mood, stress, and previous sexual experiences
    \item \textbf{Physical examination:} where indicated, including pelvic examination if pain or dryness is reported
    \item \textbf{Blood tests:} thyroid function, hormone levels (oestradiol, testosterone, prolactin), and full blood count where a physical cause is suspected
    \item \textbf{Mental health assessment:} screening for depression and anxiety
    \item \textbf{Medication review:} identifying drugs that suppress desire, such as antidepressants
    \item \textbf{Partner involvement:} where appropriate, joint discussion with the partner to understand relationship factors
\end{itemize}

\section*{Management}
\begin{itemize}
    \item \textbf{Treat underlying causes:} thyroid disease, depression, and pain conditions; change or adjust contributory medicines where possible
    \item \textbf{Psychological therapy:} cognitive behavioural therapy and sex therapy address the thoughts, fears, and relationship patterns that maintain low desire
    \item \textbf{Couple therapy:} improving communication, intimacy, and realistic expectations within the relationship
    \item \textbf{Address sexual pain and dryness:} lubricants, vaginal moisturisers, and vaginal oestrogen after menopause make sex comfortable again
    \item \textbf{Hormone treatments:} menopausal hormone therapy can improve desire in some women; testosterone therapy, used off-label and under specialist supervision, helps a minority with low desire after menopause
    \item \textbf{Flibanserin and bremelanotide:} medicines approved overseas for low desire in premenopausal women, used with caution and specialist input
    \item \textbf{Sensate focus:} structured exercises that rebuild intimacy and pleasure without the pressure of intercourse
    \item \textbf{Lifestyle:} sleep, stress reduction, exercise, and moderating alcohol
\end{itemize}

\section*{Complications}
\begin{itemize}
    \item Relationship conflict, resentment, and breakdown
    \item Reduced intimacy and emotional distance between partners
    \item Low self-esteem, anxiety, and worsening depression
    \item Avoidance of medical care if the issue is not raised
    \item Infertility of the couple if desire is low during family planning, though most couples conceive with support
\end{itemize}

\section*{Prognosis and Outcomes}
Most women with low desire improve with the right help. When a physical cause or medicine is identified, addressing it often restores desire. When the causes are psychological or relational, sex therapy and couple counselling produce good results, especially when both partners are engaged. Normalising the wide natural variation in desire is itself therapeutic: many women find that understanding their responsive desire pattern relieves distress even before any treatment. For those with persistent distressing low desire, a specialist multidisciplinary approach offers the best outcomes.

\section*{Prevention and Self-Care}
\begin{itemize}
    \item Talk openly with a partner about desire, stress, and expectations, rather than avoiding the topic
    \item Make time for intimacy and connection beyond intercourse, including touch and conversation
    \item Prioritise sleep, manage stress, and maintain regular physical activity
    \item Moderate alcohol and review recreational drugs, which suppress desire
    \item Ask about the sexual side effects of any new medicine
    \item Use lubricants and address pain early rather than enduring it
    \item See a doctor, psychologist, or sex therapist early; low desire is common and help is available
\end{itemize}

\section*{Sources and Further Reading}
The following reputable sources provide further information for healthcare students and patients seeking reliable, current advice about loss of female libido.
\begin{itemize}
    \item Healthdirect Australia. \textit{Loss of libido.} https://www.healthdirect.gov.au/loss-of-libido
    \item Jean Hailes for Women's Health. \textit{Sexual health and desire.} https://www.jeanhailes.org.au/
    \item True Relationships \& Reproductive Health. \textit{Sexual health services.} https://www.true.org.au/
    \item NHS. \textit{Loss of libido.} https://www.nhs.uk/conditions/loss-of-libido/
    \item Mayo Clinic. \textit{Female sexual dysfunction.} https://www.mayoclinic.org/
    \item MedlinePlus. \textit{Female sexual dysfunction.} https://medlineplus.gov/
\end{itemize}
